\chapter{What you can and cannot learn}
\label{ch:identif}

\begin{motivation}
A responsible fit reports not just point estimates but \emph{which} of them to
trust. This chapter explains, with the information geometry behind it, why the
branching ratio $\kappa$ is well identified from counts while the decay $\theta$ is
not; how to get standard errors that are honest under the Hawkes over-dispersion;
and how to check the fit at all (goodness-of-fit).
\end{motivation}

\section{The $\kappa$--$\theta$ ridge}

Write $\phi=\alpha e^{-\beta t}$ with $\alpha=\kappa\theta$, $\beta=\theta$. The
data inform two nearly-orthogonal things: the \emph{branching ratio}
$\kappa=\alpha/\beta=\int\phi$ (how much: the offspring fraction, visible in counts
and cross-correlations) and the \emph{timescale} $1/\theta$ (how fast: visible only
in timing).

\begin{whybox}[: counts hide the timescale]
A bucket compensator $\Xi_i=\int_{o_{i-1}}^{o_i}\xi$ \emph{integrates away} the
within-bucket timing that carries $\theta$. So $\partial\Xi_i/\partial\theta$ is
small and sign-oscillating while $\partial\Xi_i/\partial\kappa$ is order-one and
monotone: the Fisher information is ill-conditioned along the
$\alpha/\beta$-preserving ridge. As buckets coarsen, the $\theta$-information
$\to0$ while $\kappa$ persists.
\end{whybox}

\begin{intuitionbox}
Counts tell you \emph{how much} activity is self-generated ($\kappa$) but not
\emph{how fast} the echo decays ($\theta$). Treat $\kappa$ (and, later, covariate
effects) as trustworthy; fix or prior-constrain $\theta$ from domain knowledge or
finer data, and report it with wide intervals.
\end{intuitionbox}

\section{Standard errors under over-dispersion}

The observed information (Hessian of the IC-LL) gives model-based SEs. But true
Hawkes counts are \emph{over-dispersed} relative to the MBPP's Poisson assumption
(offspring clustering), so the naive SEs are too small. The package applies a
\emph{quasi-Poisson} correction --- inflate by $\sqrt{\text{dispersion}}$:

\begin{lstlisting}[caption={Standard errors with the quasi-Poisson correction (from \code{interval\_censored.py}).}]
se = _observed_information_se(objective, best.x, build_to_natural)   # Hessian^-1, delta method
disp = np.sum((counts - Xi_opt)**2 / Xi_opt) / (m - 2)               # Pearson dispersion
se_kappa, se_theta = se[0]*np.sqrt(max(disp,1.0)), se[1]*np.sqrt(max(disp,1.0))
\end{lstlisting}

\begin{examplebox}[: coverage]
On \emph{model-correct} (MBPP-Poisson) data the 95\% interval for $\kappa$ covers
the truth $\approx94\%$ of the time --- the SE computation is right. On genuine
over-dispersed Hawkes data the naive SE undercovers (57\%); the quasi-Poisson
correction lifts it to $\approx82\%$. The residual gap is serial clustering a
scalar dispersion cannot fully capture --- honest, and a reason to also run a
goodness-of-fit check.
\end{examplebox}

\section{Goodness of fit}

Two complementary checks (Appendix~\ref{app:estimation}, \code{gof.py}):

\paragraph{Time-rescaling (event-time).} If the fitted compensator is right, the
rescaled inter-event gaps $\Lambda(t_k)-\Lambda(t_{k-1})$ are i.i.d.\
$\mathrm{Exp}(1)$; a Kolmogorov--Smirnov test against $\mathrm{Exp}(1)$ detects
misspecification.

\paragraph{Dispersion (interval-censored).} The Pearson statistic
$\sum_i(\sfont C_i-\Xi_i)^2/\Xi_i$ divided by degrees of freedom should be near 1;
$>1$ flags over-dispersion (unmodelled clustering), $\gg1$ a wrong model.

\begin{lstlisting}
from hawkes_calibration import ks_test_exp1, time_rescaling_residuals, dispersion
D, p = ks_test_exp1(time_rescaling_residuals(events, compensator))   # event-time
phi  = dispersion(counts, Xi, n_params=2)                            # interval counts
\end{lstlisting}

\begin{examplebox}[: GOF separates models]
In \code{exp10\_deep.py}, the dispersion under the \emph{correct} kernel clusters
near $1.9$ (the inherent Hawkes over-dispersion), while a deliberately
\emph{wrong} kernel gives dispersion $\approx67$ --- a decisive separation.
\end{examplebox}

\begin{pitfall}
Never report $\hat\theta$ from interval-censored data as if it were as reliable as
$\hat\kappa$. And a dispersion near 1 does not by itself prove the model --- combine
it with the time-rescaling test when timestamps are available.
\end{pitfall}

\begin{recap}
$\kappa$ is well identified, $\theta$ weakly so (counts integrate away timing).
Honest SEs need a quasi-Poisson inflation for Hawkes over-dispersion; goodness-of-fit
uses time-rescaling (KS vs Exp(1)) and the Pearson dispersion. With Part II
complete --- model, solve, fit, diagnose --- Part III adds covariates.
\end{recap}

\exercise{Explain why coarser buckets make $\theta$ \emph{less} identifiable but
leave $\kappa$ essentially unchanged.}
\exercise{Simulate over-dispersed counts and verify the dispersion statistic
exceeds 1; then confirm the quasi-Poisson SE widens accordingly.}
